\section{曲線で囲まれた面積}\label{節:曲線間の面積}

$a<b$ とし，$f,g$ は $[a,b]$ で連続で，$f(x)\geq g(x)$ を満たすとする。二曲線 $y=f(x)$，$y=g(x)$ の間を，$x$ 軸に垂直な細い帯に分ける。幅が $\Delta x$ の帯の高さは，その付近での $f(x)-g(x)$ である。したがって面積は
\begin{equation}\label{式:上下の差の面積}
S=\int_a^b\{f(x)-g(x)\}\dd x
\end{equation}
になる。上下の曲線がともに $x$ 軸の下にあっても，高さは上の値から下の値を引いた差であり，同じ式が使える。

二曲線が途中で交差するなら，上下関係が変わる点で区間を分ける。一つの式で表せば
\[
S=\int_a^b|f(x)-g(x)|\dd x
\]
であるが，実際に計算するには差の符号を調べて絶対値を外す。交点の $x$ 座標は区切りの候補であり，交点を求めただけで上下が決まるわけではない。差を因数分解したり，各区間で一つの値を代入したりして確かめる。

\begin{example}{途中で上下が入れ替わる面積}
二曲線 $y=x^3-3x$ と $y=x$ に囲まれた部分の面積を求めよ。
\end{example}

交点は $x^3-3x=x$，すなわち $x(x-2)(x+2)=0$ より，$x=-2,0,2$ である。差 $x^3-4x$ は $-2<x<0$ で正，$0<x<2$ で負なので，面積は
\[
S=\int_{-2}^0(x^3-4x)\dd x
 +\int_0^2(4x-x^3)\dd x.
\]
図\ref{図:面積の上下交換}では，左右の部分が原点について対称である。積分でも第一項に $x=-t$ を使えば第二項と同じになるから，
\[
S=2\int_0^2(4x-x^3)\dd x
 =2\left[2x^2-\frac{x^4}{4}\right]_0^2=8.
\]
これを $\int_{-2}^2(x^3-4x)\dd x$ とすれば奇関数の積分なので零になるが，その値は左右の面積を一方は正，一方は負として足したものである。

\begin{figure}[htbp]
\centering
\begin{tikzpicture}[x=1.6cm,y=0.65cm,>=Stealth]
\fill[blue!12]plot[domain=-2:0,samples=60](\x,{\x*\x*\x-3*\x})--plot[domain=0:-2,samples=30](\x,\x)--cycle;
\fill[orange!17]plot[domain=0:2,samples=30](\x,\x)--plot[domain=2:0,samples=60](\x,{\x*\x*\x-3*\x})--cycle;
\draw[->](-2.3,0)--(2.4,0)node[right]{$x$};
\draw[->](0,-2.8)--(0,3.4)node[above]{$y$};
\draw[blue!70!black,thick]plot[domain=-2.13:2.13,samples=120](\x,{\x*\x*\x-3*\x});
\draw[orange!80!black,thick](-2.2,-2.2)--(2.2,2.2);
\foreach \x in {-2,2}{\draw[dashed](\x,0)--(\x,\x);\node[below]at(\x,0){$\x$};}
\node[below right]at(0,0){$O$};
\node[blue!70!black]at(-1.25,2.8){$y=x^3-3x$};
\node[orange!80!black]at(1.6,3.0){$y=x$};
\end{tikzpicture}
\caption{同じ二曲線でも，交点の左右で上側の曲線が入れ替わる。}\label{図:面積の上下交換}
\end{figure}

\subsection{横に切るほうが簡単な場合}

細い帯は縦に取る必要はない。$c\leq y\leq d$ において，右側の境界が $x=q(y)$，左側が $x=p(y)$ と表せるなら，帯の横の長さは $q(y)-p(y)$ なので
\begin{equation}\label{式:左右の差の面積}
S=\int_c^d\{q(y)-p(y)\}\dd y
\end{equation}
である。式\eqref{式:上下の差の面積}では上下の差，この式では左右の差を使う。変数の文字を交換するだけでなく，実際にどの向きに帯を取るかが変わっている。

たとえば $x=y^2$ と $x=2-y$ に囲まれる部分では，交点の $y$ 座標は $-2,1$ である。その間では直線が右側なので，
\[
S=\int_{-2}^1(2-y-y^2)\dd y
 =\left[2y-\frac{y^2}{2}-\frac{y^3}{3}\right]_{-2}^1
 =\frac92.
\]
縦に切ると放物線の上側と下側を $y=\pm\sqrt x$ に分け，さらに直線との関係も場合分けすることになる。横に切れば，一つの区間で一つの差になる。境界をどちらの変数で表すと少ない場合分けで済むかが，積分変数を選ぶ基準になる。

%
